% generated by figures/make_figures.py — do not edit
\begin{table}[H]
\caption{Recorded etymologies on Wikipedia by lexicon class and by language of the name (post hoc, amendment 5). Upper part: among the names in each class of the lexicon whose Wikipedia section states a derivation, the share whose recorded derivation is in the same class and the share whose derivation is water-related in any class (the highland class is not water-related). Lower part, for languages with at least 30 coded sections: the share of water-related derivations among the names that the lexicon read as water-related (its precision) and among the other names, with counts. Shares are unweighted; intervals are 95\,\% \citep{wilson1927}.}\label{tab:etym-detail}
\footnotesize
\resizebox{\ifdim\width>\textwidth\textwidth\else\width\fi}{!}{%
\begin{tabular}{lccc}
\tophline
Lexicon class & Coded & Same class & Water-related, any class \\
\middlehline
water & 38 & 82\,\% (67--91) & 84\,\% (70--93) \\
wetland & 6 & 50\,\% (19--81) & 100\,\% (61--100) \\
riverine & 11 & 73\,\% (43--90) & 82\,\% (52--95) \\
crossing & 18 & 89\,\% (67--97) & 94\,\% (74--99) \\
lowland & 37 & 49\,\% (33--64) & 54\,\% (38--69) \\
highland & 15 & 53\,\% (30--75) & 27\,\% (11--52) \\
\middlehline
Language of the name & Coded & Water-related, read as water-related & Water-related, others \\
\middlehline
English & 48 & 75\,\% (55--88); 18/24 & 29\,\% (15--49); 7/24 \\
French & 120 & 67\,\% (50--80); 22/33 & 18\,\% (12--28); 16/87 \\
Polish & 35 & 40\,\% (12--77); 2/5 & 3\,\% (1--17); 1/30 \\
Slovene & 54 & 94\,\% (73--99); 16/17 & 16\,\% (8--31); 6/37 \\
\bottomhline
\end{tabular}}
\end{table}
